% Transcription — fonds Alexandre Grothendieck, shelfmark 70, batch 7 (pages 121-125).
% Established from the facsimile archives/batches/70/batch-07.pdf (a mirror of
% https://grothendieck.umontpellier.fr/70.pdf), per the skill
% transcribe-grothendieck.
%
% Pass: Opus 5.5 (claude-opus-5-5), 2026-09-29 — first pass, unchecked against the pages by a human.
%
% The last, short batch of the folder: five pages, three transcribed.
%   Pages 121-122 — "Programme provisoire d'une théorie arithmétique des
%     polyèdres réguliers": the canonical representative X_tau of a type tau,
%     its ordered set of faces and its group G_tau; then the definition of a
%     k-realisation (projective space P over k with s_0, H_infty, and a set of
%     projective subspaces subject to conditions 1)-4)).
%   Page 123 — a sheet of notation (s_i, s'_i, delta_i, f_i, N = card G, the
%     generators sigma_i) with sketches of a flag in a cube and a polygon.
%   Page 124 is an unrelated typed verso (a USTL letter stamped 17 JUN 1976 —
%     consistent with the inventory's "[à partir de 1976]") and page 125 a blank
%     coloured sheet; both are skipped.
%
% Notation fixed for the batch: his looped script capital for the set of faces
% is set \mathcal{F} (see the note at page 121).

\documentclass[11pt,a4paper]{article}
\input{../preamble/grothendieck}

\folder{70}
\batch{7}
\pages{121}{125}
\dating{[à partir de 1976]}
\watermark{Édition de démonstration}
\foldertitle{Réalisations géométriques de structures combinatoires (n-polyèdres, n-hyperpolyèdres [2-polyèdres réguliers]…) (Vieilles rédactions) : notes manuscrites (s.d.).}

\begin{document}

\section{Programme provisoire d'une théorie arithmétique des polyèdres réguliers}

\page{121}
\note{Le titre de la section est le sien : il ouvre la page, souligné. La page commence un texte neuf ; rien n'y reprend une page antérieure.}

Soit un type \add{$\tau$} de polyèdre régulier, dans l'esp. affine de dim $n$
(sur $\mathbb{R}$). Il y a un représentant $X_\tau$ « canonique » de $\tau$,
défini à \uncertain{isom. can.} près, par le choix d'un « repère » $R_\tau$ de
$X_\tau$. L'espace affine $E_\tau$ qui le porte est muni d'une origine (le
\struck{centre} \add{bary}centre \struck{de gravité} de l'ens. des pts
extrémaux), donc c'est un vectoriel.

On désigne par $\mathcal{F}_\tau$ l'ensemble ordonné des facettes (y compris
celles de dim $-1$ et $n$) de $X_\tau$, par $G_\tau$ le groupe
\[
  G_\tau \simeq \mathrm{Aut}_{\mathrm{aff}}(X_\tau) \simeq \mathrm{Aut}_{\mathrm{ord}}(\mathcal{F}_\tau)
\]
\note{La lettre qu'il emploie pour l'ensemble des facettes est une capitale cursive bouclée, entre $\Phi$ et un $\mathcal{F}$ ; on la rend partout par $\mathcal{F}$. L'indice du second $\mathrm{Aut}$ se lit « ord » ou « ad » ; « ord » est retenu d'après le contexte.}

La catégorie des polyèdres réguliers de type $\tau$ équivaut à celle des
torseurs sous $G_\tau$.

Soit $k$ un anneau commutatif (plus gén., on pourrait se placer sur un schéma,
ou sur un topos loc. annelé -- ici on est sur $\operatorname{Spec} k$).
Une $k$-réalisation (\struck{\uncertain{de plongée}}
\add{\uncertain{déf. ci-dessous}}) de \struck{$X$} $\tau$ (ou de $X_\tau$)
est par définition la donnée de
\marginal{raisonnable si $\operatorname{Spec} k$ connexe…}

a) Un $k$-espace projectif \add{$P$} de dim $n$ (splitté, donc associé à un
module projectif de rg.~f. défini \ill{} tensorisation près par un
\struck{\ill{}} module inversible) [muni d'une section $s_0$ et d'un hyperplan
(« à l'infini ») $H_\infty$]

b) Un ensemble $\mathcal{F}$ de sous-espaces projectifs \struck{$P$} de $P$,
\uncertain{donnés} soumis aux conditions suivantes :

\page{122}
1) $\mathcal{F}$, pour la relation d'ordre induite par l'inclusion, est
isomorphe à $\mathcal{F}_\tau$.

[\struck{2)} \struck{Pour} Pour tout $V \in \mathcal{F}$, on a
$s_0 \notin V$ et $V \not\subset H_\infty$ \emph{fibre par fibre}]
\note{Le passage entre crochets porte un « 2) » noirci ; une flèche le contourne et rattache la suite directement à 1).}

et pour $V \in \mathcal{F}$, la dimension combinatoire de
\uncertain{$\mathcal{F}$} dans l'ens. ordonné $\mathcal{F}$ (en prenant $-1$
comme dim du plus petit élt) est \emph{égale} à la dim. projective de
$V$. \struck{\ill{}} Enfin $P \in \mathcal{F}$ (donc $\dim P = n$)
\note{« la dimension combinatoire de $\mathcal{F}$ » : on attendrait $V$ ; la page porte la lettre de l'ensemble.}

3) \struck{Tout} Les éléments de $\mathcal{F}_0$ engendrent $P$
projectivement (fibre par fibre), et ceux de $\mathcal{F}_{n-1}$ engendrent
$\check{P}$. \struck{(Ceci implique qu'un automorphisme de $P$ qui
\uncertain{induise} $\mathcal{F}$ est connu quand on connaît sa restriction à
$\mathcal{F}$ comme \struck{automorphisme} \add{permutation} de l'ens. ordonné
$\mathcal{F}$…)}
\note{La parenthèse est encadrée et barrée de traits obliques ; un trait la relie à la note marginale qui suit, qui la remplace.}

\marginal{Plus précisément, on veut que pour tout $V \subset W$,
$V, W \in \mathcal{F}$, $W \neq V$, soit égal le sous-espace projectif de $P$
engendré par les $V' \in \mathcal{F}$ tels que $V \subset V' \subset W$,
$\dim V' = \dim V + 1$ -- et \ill{} dual}

4) Tout automorphisme de l'ens. ordonné $\mathcal{F}$ est induit par un
automorph. de $P$ invariant $\mathcal{F}$.

\page{123}
\note{Feuille de notations, sans texte suivi ; les symboles $s_i$, $s'_i$, $\delta_i$, $f_i$ ne sont définis sur aucune page de ce lot.}

\[
\begin{array}{lll}
  (s_i)_{0 \leq i \leq n-1} & s_{n-1} = 2 & s_n = 1 \\
  (s'_i)_{0 \leq i \leq n-1} & s'_0 = 2 & s'_{-1} = 1 \\
  (\delta_i)_{0 \leq i \leq n-2} & & \\
  (f_i)_{0 \leq i \leq n-1} & f_{-1} = f_n = 1 &
\end{array}
\]
\[
  N = \prod_{1}^{n-1} s_i = \prod_{1}^{n-1} s'_i = \operatorname{card} G
\]
\[
  \begin{cases}
    s'_i f_{i+1} = f_i s_{i+1} \quad (0 \leq i \leq n-1) \\
    s
  \end{cases}
  \qquad
  \frac{f_i}{s'_i} = \frac{f_{i+1}}{s_{i+1}} = f
\]
\note{La seconde ligne de l'accolade s'arrête sur « $s$ ». Au numérateur de la fraction, l'indice de $f$ est peu net.}

\[
  x_{i-1} \;\bullet\; \bullet\; x_{i+2} \qquad \sigma_i \sigma_{i+1}
\]
\[
  \sigma_0\sigma_1,\ \sigma_1\sigma_2,\ \dots,\ \sigma_{n-2}\sigma_{n-1}
  \qquad G \ni \sigma_0, \dots, \sigma_{n-1}
\]
\note{Le symbole entre $G$ et $\sigma_0$ peut aussi se lire $\supset$.}

\drawing{À droite, un polygone régulier (hexagone ou heptagone) dont un sommet et les deux arêtes qui en partent sont repassés en gras, avec un rayon vers le centre ; au-dessous, un cube en perspective où un drapeau (sommet, arête, face) est marqué en gras, avec trois droites étiquetées $\sigma_0$, $\sigma_1$, $\sigma_2$. À côté, la chaîne}

\[
\begin{array}{cccccccc}
  x_0 & x_1 & & x_{i-1} & x_i & & & x_{n-1} \\
  F_0 & F_1 & - & F_{i-1} & F_i & F_{i+1} & - & F_{n-1}
\end{array}
\]
\note{Une accolade sous $F_0,\dots,F_{i-1}$ est marquée $V_{i-1}$, une autre sous $F_{i+1},\dots,F_{n-1}$ est marquée $V_i$. La lettre des $F$ est la même capitale cursive que ci-dessus ; le premier est surchargé par le $\sigma_2$ du dessin.}

\drawing{Plus bas, « $x_0$ --- $x_{n-1}$ » au-dessus d'un segment gradué fléché, dont les deux premiers intervalles sont repassés en gras.}

\end{document}
